%!TEX root = /Users/markelikalderon/Documents/Git/perceptual_self-consciousness/perceptual_self-consciousness.tex
\chapter{Introduction} % (fold)
\label{cha:introduction}

\section{On the Semantic Field of ``Consciousness''} % (fold)
\label{sec:semantic_field}

The present essay is about perceptual self-consciousness. I shall approach my topic through the consideration of an ancient \emph{aporia}, first put forward by Plato in the \emph{Charmides} (chapter~\ref{cha:offering}), developed by Aristotle in \emph{De anima} 3.2 (chapter~\ref{cha:perceiving}), and commented upon by the Neoplatonists (chapter~\ref{cha:procession}).

But first, what is meant by perceptual self-consciousness? After all, consciousness is spoken of in many ways, as Aristotle might have put it (\citealt[2.2.1]{Brentano:1874aa}, \citealt{lycan87}). Let us begin, \emph{via negativa}, by considering what is not meant.

In a fascinating study of the history and afterlife of Aristotle's common power (\emph{koinē dunamis}) or common sense (\emph{koinē aisthēsis}), \citet[21]{Heller-Roazen:2007ee} claims that ``consciousness'' is ``the name for the cognitive faculty by which we represent our thoughts to ourselves.'' If we accept Heller-Roazen's conception of consciousness and accept, as well, the perception–cognition divide, on some sutiable understanding of that distinction, then talk of perceptual self-consciousness would be an oxymoron. Indeed, it is on this basis that Heller-Roazen criticizes modern commentators for deploying the term in their exegesis of ancient texts that discuss perceptual powers. 

It is true that there is no term in Greek or Latin with the full range of senses conveyed by the modern word ``consciousness'' (\citealt[21, 302 n2]{Heller-Roazen:2007ee}, for a history see \citealt[chapter 8]{Lewis:1960ga})—not even the Latin \emph{conscientia}, which can mean privy, to know along with others, and assumed a moral sense in a proleptic anticipation of the modern conception of a (good or bad) conscience. Indeed, the English word ``consciousness'' is unrelated to the Latin \emph{conscientia}, being first introduced by Ralph Cudworth in \emph{The True Intellectual System of the Universe} (1678, 159) as a translation of Plotinus's use of the Greek word \emph{sunaisthēsis}. This is not, indeed, a bad Latinate transliteration, though \emph{aisthēsis} may, though need not (chapter~\ref{sub:a_schematic_representation}), carry perceptual connotations not shared with \emph{scientia} (on Cudworth on consciousness see \citealt{Thiel:1991hw}, \citealt[chapter 2]{Thiel:2011cq}, and \citealt{Leisinger:2022sh}). Be that as it may, I hope to show that the specific psychic phenomenon that I am interested was in fact discussed in antiquity and that we can learn from that discussion.

So by ``consciousness'', I do not mean a cognitive faculty by which we represent our thoughts to ourselves. Consciousness is not the sole or exclusive provenance of reason. Rather, I have in mind a distinctively perceptual phenomenon. 

Neither do I have in mind phenomenal consciousness, at least if the qualifier is meant to pick out the qualitative character of experience as discussed by, \emph{inter alia}, \citet{Jackson:1982my}. (The phenomenal too is spoken of in many ways. Originally it is the character of what appears in experience, and only subsequently the character of the appearance itself.) 

In a section ominously entitled \emph{The Snares of Consciousness}, McCabe \citeyearpar{McCabe:2007jb} puruses the charge of anachronism against \citet{Kosman:2014ab}. Kosman, had, by McCabe's lights, misattributed to Aristotle an account of phenomenal consciousness but no such conception is alleged to be found in the Aristotelian corpus. This, according to McCabe, is a modern intrusion that obscures more than it enlightens.

The terms of the charge are surprising, however, since \citet{Kosman:2014ab} never speaks of phenomenal consciousness. He speaks instead of reflexive self-awareness or self-consciousness. And in an earlier, related, work on the \emph{Charmides}, \citet{Kosman:2014aa} speaks as well of self-presence. Far from being a modern imposition upon ancient texts, this latter notion traces back to late anitiquity (Porphyry \emph{Sententiae} 40–1, Augustine \emph{De Trinitate}, perhaps not unrelated if Porphyry's work was among the books of the Platonists that inaugurates Augustine's philosophical conversion, \emph{Confessiones} 7.9.13, for discussion see \citealt[chapter 12]{Sorabji:2006aa}). Regardless of the justice of McCabe's charge of anachronism, it is irrelevant to the present study, since I am not at all concerned with the qualitative character of perceptual experience. 

So by ``consciousness'', I do not mean phenomenal consciousness, at least if that is meant to designate the qualitative character of experience. Rather, I have in mind a kind of perceptual self-awareness—being aware from within that one perceives what one does. 

Some comment is required about the prefix ``self-'' in ``perceptual self-conscious''.

In English, there is a sense of ``self-conscious'', familiar from ordinary life, that emerged in the early 19th century. Consider Pip's first tragic encounter with Estella:
\begin{quotation}
	``He calls the knaves, Jacks, this boy!'' said Estella with disdain, before our first game was out. ``And what coarse hands he has! And what thick boots!'' 

	I had never thought of being ashamed of my hands before; but I began to consider them a very indifferent pair. Her contempt for me was so strong, that it became infectious, and I caught it. (Charles Dickens, \emph{Great Expectations}, chapter 8)
\end{quotation}
Estella's disdain for Pip's coarse hands makes him self-conscious about his hands in way that he previously had not been. In this sense, to be self-conscious is to be excessively or unduly aware of one's person or some aspect of one's appearance, often, like Pip, prompted by social circumstance and often, again like Pip, causing shame or some other form of distress. Being self-conscious, in this sense, has a tendency to make one less spontatneously active, less authentic, and more anxious. While not without philosophical interest, given its obvious connections with \emph{mauvaise foi}, this is not a sense of self-consciousness that I shall be discussing.

There is, however, a notion of self-consciousness that I shall discuss, namely, reflective self-consciousness, if only to distinguish it from my proper topic. By ``reflective self-consciousness'', I mean an awareness of our perceptual experience afforded by reflection, understood as a distinct mental act that takes the perceptual experience as its object. By contrast, my proper topic is a mode of awareness immanent in the perceptual act itself, in the way that the awareness afforded by reflection is manifestly not. Such awareness is better described as reflexive than reflective.

% section semantic_field (end)

\section{Non-Positional Consciousness} % (fold)
\label{sartre}

We can, sometimes at least, be aware from within that we perceive what we do. Moreover, we can be aware from within that we perceive what we do while we do. The awareness of so perceiving unfolds along with the perceiving. If the relevant mode of awareness is immanent in the perceptual act, it is perhaps unsurprising that they should share a temporal aspect.

Suppose we accept that perceptual self-consciousness reflexive rather than reflective. There are two ways to understand how reflexive self-awareness may be immanent to the perceptual act.

First, being aware from within of perceiving what one perceives may be understood as a form of joint perception. Estella in perceiving Pip's coarse hands is aware that she does. Perhaps she jointly perceives Pip's coarse hands and her seeing of those hands. \citet{Twomey:2019aa} attributes such a view to Aristotle on the basis of the opening \emph{aporia} of \emph{De anima} 3.2. If the perception jointly perceives Pip's hands and Estella's seeing of them, then there is a straightforward sense in which reflexive awareness is immanent in Estella's perceptual activity. If self-aware perception is a joint perception, then not only is Pip's hands the intentional object of the perceptual act but so is the perceptual act itself. 

The second way of understanding immanent self-awareness denies that the perception is itself an intentional object. Sartre in a number of early writings embraces this denial. For Sartre, Estella's consciousness of Pip's coarse hands is positional (\emph{positionelle})
\begin{quote}
	All consciousness is positional in that it transcends itself to reach an object, and it is exhausted by just this act of positing. (Sartre, \emph{L\^{E}tre et le N\'{e}ant}; Richmond, \citeyear[10]{Sartre:2018aa})
\end{quote}
In seeing Pip's coarse hands, Estella's consciousness transcends itself to to reach them. What makes talk of transcendence, here, apt is that the object of Estella's vision is other than the conscious act directed upon it. Sartre maintains that the intentional object of a conscious act is other that the conscious act directed upon it which is why in order to reach its object that act must transcend itself.

This has an important consequence for the reflexive case.

% section sartre (end)

\section{Looking Forward} % (fold)


% section looking_forward (end)

% chapter second_excurus (end)